% =====================================================================
%  SUPERSEDED 2026-08-19 -- NOT PART OF THE BUILD.
%  Figure 1 is now figures/fig1_objective.tex, wrapped by
%  figures/fig1_objective_float.tex, which sections/01_intro.tex \input's.
%  Nothing \input's THIS file, so nothing in it renders.
%
%  Its notation is deliberately left as \log p_\theta rather than being
%  migrated to \log q_\theta: the file also asserts the retired equality
%  p_\theta(.|c) = pi(.|c) (line ~309), a claim the paper no longer makes.
%  Migrating the symbol alone would leave a retired claim wearing current
%  notation.  If this figure is ever revived, rewrite the claim first.
% =====================================================================
% =====================================================================
%  fig1_method.tex --- BGFM / HBFM Figure 1 (method overview)
%
%  Pure TikZ.  No external images, no pgfplots, no \includegraphics,
%  no arrows.meta (TeX Live 2018 compatible).
%  Okabe--Ito colour-blind-safe palette (deuteranopia/protanopia safe).
%
%  USAGE
%   (a) standalone (verification / slides):   pdflatex fig1_method.tex
%   (b) inside the paper -- main.tex must \usepackage{tikz} and
%       \usepackage{amsmath} in its preamble:
%         \begin{figure}[t]
%           \centering
%           % =====================================================================
%  SUPERSEDED 2026-08-19 -- NOT PART OF THE BUILD.
%  Figure 1 is now figures/fig1_objective.tex, wrapped by
%  figures/fig1_objective_float.tex, which sections/01_intro.tex \input's.
%  Nothing \input's THIS file, so nothing in it renders.
%
%  Its notation is deliberately left as \log p_\theta rather than being
%  migrated to \log q_\theta: the file also asserts the retired equality
%  p_\theta(.|c) = pi(.|c) (line ~309), a claim the paper no longer makes.
%  Migrating the symbol alone would leave a retired claim wearing current
%  notation.  If this figure is ever revived, rewrite the claim first.
% =====================================================================
% =====================================================================
%  fig1_method.tex --- BGFM / HBFM Figure 1 (method overview)
%
%  Pure TikZ.  No external images, no pgfplots, no \includegraphics,
%  no arrows.meta (TeX Live 2018 compatible).
%  Okabe--Ito colour-blind-safe palette (deuteranopia/protanopia safe).
%
%  USAGE
%   (a) standalone (verification / slides):   pdflatex fig1_method.tex
%   (b) inside the paper -- main.tex must \usepackage{tikz} and
%       \usepackage{amsmath} in its preamble:
%         \begin{figure}[t]
%           \centering
%           % =====================================================================
%  SUPERSEDED 2026-08-19 -- NOT PART OF THE BUILD.
%  Figure 1 is now figures/fig1_objective.tex, wrapped by
%  figures/fig1_objective_float.tex, which sections/01_intro.tex \input's.
%  Nothing \input's THIS file, so nothing in it renders.
%
%  Its notation is deliberately left as \log p_\theta rather than being
%  migrated to \log q_\theta: the file also asserts the retired equality
%  p_\theta(.|c) = pi(.|c) (line ~309), a claim the paper no longer makes.
%  Migrating the symbol alone would leave a retired claim wearing current
%  notation.  If this figure is ever revived, rewrite the claim first.
% =====================================================================
% =====================================================================
%  fig1_method.tex --- BGFM / HBFM Figure 1 (method overview)
%
%  Pure TikZ.  No external images, no pgfplots, no \includegraphics,
%  no arrows.meta (TeX Live 2018 compatible).
%  Okabe--Ito colour-blind-safe palette (deuteranopia/protanopia safe).
%
%  USAGE
%   (a) standalone (verification / slides):   pdflatex fig1_method.tex
%   (b) inside the paper -- main.tex must \usepackage{tikz} and
%       \usepackage{amsmath} in its preamble:
%         \begin{figure}[t]
%           \centering
%           % =====================================================================
%  SUPERSEDED 2026-08-19 -- NOT PART OF THE BUILD.
%  Figure 1 is now figures/fig1_objective.tex, wrapped by
%  figures/fig1_objective_float.tex, which sections/01_intro.tex \input's.
%  Nothing \input's THIS file, so nothing in it renders.
%
%  Its notation is deliberately left as \log p_\theta rather than being
%  migrated to \log q_\theta: the file also asserts the retired equality
%  p_\theta(.|c) = pi(.|c) (line ~309), a claim the paper no longer makes.
%  Migrating the symbol alone would leave a retired claim wearing current
%  notation.  If this figure is ever revived, rewrite the claim first.
% =====================================================================
% =====================================================================
%  fig1_method.tex --- BGFM / HBFM Figure 1 (method overview)
%
%  Pure TikZ.  No external images, no pgfplots, no \includegraphics,
%  no arrows.meta (TeX Live 2018 compatible).
%  Okabe--Ito colour-blind-safe palette (deuteranopia/protanopia safe).
%
%  USAGE
%   (a) standalone (verification / slides):   pdflatex fig1_method.tex
%   (b) inside the paper -- main.tex must \usepackage{tikz} and
%       \usepackage{amsmath} in its preamble:
%         \begin{figure}[t]
%           \centering
%           \input{figures/fig1_method}
%           \caption{...}\label{fig:method}
%         \end{figure}
%
%  SIZE.  Natural size 13.95cm x 8.70cm == 5.49in x 3.43in.
%  ICLR is a SINGLE 5.5in-wide text column (10pt Times, 11pt leading), so
%  this is a full-text-width figure at scale 1.0.  DO NOT rescale it: every
%  glyph inside is \footnotesize / \scriptsize / \tiny of the 10pt body
%  font, and scaling desynchronises in-figure text from the caption.
%  If a shorter figure is needed, delete band (e) (the total-loss box and
%  legend, y < 1.7) and state the loss in the caption instead -- that saves
%  1.5cm without touching any other coordinate.
%
%  LAYOUT (five bands, top to bottom).  Wire routing is deliberately
%  crossing-free:
%    (a) generative spine   prior noise -> flow matching -> x_1
%    (b) global-shape term  FFJORD log p  vs  E/kT     -> L_energy  [LEFT]
%    (c) local-gradient term closed-form score vs F/kT -> L_force   [RIGHT]
%    (d) physics bus        universal neural potential (OMol25 / eSEN)
%    (e) total loss + legend
%  (b) is on the left because the reverse-time (FFJORD) arrow terminates at
%  t=0 on the left; (c) is on the right because the score is read out at
%  t*=0.95 on the right; (d) is a full-width bus underneath so both terms
%  are fed by short vertical arrows and no wire ever crosses another.
% =====================================================================

% ---- standalone bootstrap (skipped when \input into a tikz-enabled doc) ----
\ifx\tikzpicture\undefined
  \documentclass[10pt]{article}
  \usepackage[paperwidth=14.30cm,paperheight=9.15cm,margin=0.15cm,noheadfoot]{geometry}
  \usepackage{amsmath,amssymb,bm}
  \usepackage{tikz}
  \setlength{\topskip}{0pt}
  \setlength{\parindent}{0pt}
  \pagestyle{empty}
  \begin{document}%
  \def\bgfmFinish{\end{document}}
\fi
\providecommand{\bgfmFinish}{\relax}
% \pending{...}: red placeholder marker, shared with the main text so that a
% single grep finds every not-yet-real number in the manuscript.
\providecommand{\pending}[1]{\textcolor{red}{#1}}
\providecommand{\bgfmFigScale}{1.0}

% 6pt/6.6pt: one step above \tiny (5pt at 10pt base), which is below the
% conventional 6pt floor for in-figure text.
\providecommand{\bgfmtiny}{\fontsize{6}{6.6}\selectfont}
\usetikzlibrary{positioning,calc,fit,backgrounds,decorations.pathmorphing,%
                arrows,shapes.geometric}

% ---------------------------------------------------------------- palette
% Okabe & Ito (2008): the de-facto standard CVD-safe qualitative set.
\definecolor{oiOrange}{HTML}{E69F00}
\definecolor{oiSky}   {HTML}{56B4E9}
\definecolor{oiGreen} {HTML}{009E73}
\definecolor{oiYellow}{HTML}{F0E442}
\definecolor{oiBlue}  {HTML}{0072B2}
\definecolor{oiVerm}  {HTML}{D55E00}
\definecolor{oiPurple}{HTML}{CC79A7}
\definecolor{oiGrey}  {HTML}{7F7F7F}
\definecolor{oiGreyL} {HTML}{BFBFBF}
\definecolor{oiGreyD} {HTML}{555555}
%
% SEMANTIC assignment -- identical in this figure and in the matplotlib
% figures (make_experiment_figures.py::PALETTE), so the paper reads as one
% visual system:
%   oiBlue    model / flow matching / learned quantities
%   oiGreen   L_energy  (global shape, FFJORD, Boltzmann)
%   oiVerm    L_force   (local gradient, score--force matching)
%   oiOrange  physics oracle (universal neural potential, DFT labels)
%   oiPurple  negative control (Y-scrambled)
%   oiGrey    no-physics baseline / neutral structure

% ---------------------------------------------------------------- styles
\tikzset{
  bgfmpanel/.style={rounded corners=2.2pt, draw=#1, line width=0.65pt,
                    fill=#1!5},
  ttl/.style={anchor=west, font=\footnotesize\bfseries, inner sep=0pt},
  bdy/.style={anchor=west, font=\scriptsize, inner sep=0pt, align=left},
  tny/.style={anchor=west, font=\bgfmtiny, inner sep=0pt, align=left,
              text=oiGreyD},
  flowarr/.style={-stealth', line width=0.9pt, draw=oiBlue},
  revarr/.style={-stealth', line width=0.75pt, draw=oiGreen,
                 dash pattern=on 2.2pt off 1.6pt},
  physarr/.style={-stealth', line width=0.8pt, draw=oiOrange!85!black},
  tapearr/.style={-stealth', line width=0.8pt, draw=oiBlue},
  lossbadge/.style={rounded corners=2pt, draw=#1, line width=0.6pt,
                    fill=#1!12, inner sep=2.6pt, align=left,
                    font=\scriptsize},
}

% ------------------------------------------------------- molecule glyphs
% Deterministic coordinates (no \pgfmathrandom) so the figure is
% byte-reproducible across compiles and across TeX distributions.
\newcommand{\bgfmAtom}[3]{\fill[#3] (#1,#2) circle (1.45pt);%
  \draw[oiGrey!60, line width=0.15pt] (#1,#2) circle (1.45pt);}
\newcommand{\bgfmAtomS}[3]{\fill[#3] (#1,#2) circle (0.95pt);%
  \draw[oiGrey!60, line width=0.15pt] (#1,#2) circle (0.95pt);}
\newcommand{\bgfmBond}[4]{\draw[oiGrey!85, line width=0.55pt] (#1,#2) -- (#3,#4);}

% t ~ 0 : unstructured point cloud, no bonds
\newcommand{\bgfmMolNoise}{%
  \bgfmAtomS{-0.34}{ 0.30}{oiGreyL}\bgfmAtomS{ 0.10}{ 0.36}{oiGreyL}%
  \bgfmAtomS{ 0.36}{ 0.06}{oiGreyL}\bgfmAtomS{-0.10}{-0.02}{oiGreyL}%
  \bgfmAtomS{-0.40}{-0.26}{oiGreyL}\bgfmAtomS{ 0.06}{-0.34}{oiGreyL}%
  \bgfmAtomS{ 0.40}{-0.32}{oiGreyL}\bgfmAtomS{ 0.24}{ 0.22}{oiGreyL}%
  \bgfmAtomS{-0.18}{ 0.12}{oiGreyL}}

% intermediate t : partially condensed, proto-structure
\newcommand{\bgfmMolMid}{%
  \bgfmBond{-0.20}{0.14}{0.04}{0.22}\bgfmBond{0.04}{0.22}{0.24}{0.02}%
  \bgfmBond{-0.20}{0.14}{-0.14}{-0.16}%
  \bgfmAtomS{-0.20}{ 0.14}{oiGrey!70}\bgfmAtomS{ 0.04}{ 0.22}{oiGrey!70}%
  \bgfmAtomS{ 0.24}{ 0.02}{oiGrey!70}\bgfmAtomS{-0.14}{-0.16}{oiGrey!70}%
  \bgfmAtomS{ 0.30}{ 0.30}{oiGreyL}\bgfmAtomS{-0.36}{-0.24}{oiGreyL}%
  \bgfmAtomS{ 0.20}{-0.30}{oiGreyL}}

% t = 1 : a clean structure -- 5-ring + N + O + a transition metal (green),
% so the "83 elements / transition metals" claim is visible in the glyph.
\newcommand{\bgfmMolFinal}{%
  \bgfmBond{0.000}{0.280}{-0.266}{0.087}\bgfmBond{-0.266}{0.087}{-0.165}{-0.227}%
  \bgfmBond{-0.165}{-0.227}{0.165}{-0.227}\bgfmBond{0.165}{-0.227}{0.266}{0.087}%
  \bgfmBond{0.266}{0.087}{0.000}{0.280}%
  \bgfmBond{0.266}{0.087}{0.560}{0.250}%
  \bgfmBond{-0.266}{0.087}{-0.520}{0.300}%
  \bgfmBond{-0.165}{-0.227}{-0.360}{-0.440}%
  \bgfmAtom{ 0.000}{ 0.280}{oiGrey!85}\bgfmAtom{-0.266}{ 0.087}{oiGrey!85}%
  \bgfmAtom{-0.165}{-0.227}{oiGrey!85}\bgfmAtom{ 0.165}{-0.227}{oiGrey!85}%
  \bgfmAtom{ 0.266}{ 0.087}{oiBlue}%
  \fill[oiGreen] (0.560,0.250) circle (2.4pt);%
  \draw[oiGrey!70, line width=0.2pt] (0.560,0.250) circle (2.4pt);%
  \bgfmAtomS{-0.520}{ 0.300}{oiVerm}%
  \bgfmAtomS{-0.360}{-0.440}{white}}

% =====================================================================
\begin{tikzpicture}[x=1cm, y=1cm, scale=\bgfmFigScale, >=stealth']
% guarantee a stable bounding box
\path (0,0.06) rectangle (13.95,8.66);

% ---- shared geometry -------------------------------------------------
\def\Ay{6.42}\def\AY{8.52}          % band (a) vertical extent
\def\bandL{2.30}\def\bandR{10.34}   % flow band horizontal extent
\def\tz{2.55}\def\to{10.09}         % x of t=0 and t=1 on the time axis
\def\tstar{9.71}                    % x of t* = 0.95
\def\By{3.36}\def\BY{5.96}          % bands (b)/(c) vertical extent

% ==================== BAND (a): generative spine =====================
% -- prior --------------------------------------------------------------
\draw[bgfmpanel=oiBlue] (0.04,\Ay) rectangle (2.02,\AY);
\node[anchor=north west, font=\scriptsize\bfseries] at (0.12,8.48) {(a)};
\begin{scope}[shift={(1.06,7.88)}]\bgfmMolNoise\end{scope}
\node[anchor=north, font=\scriptsize] at (1.03,7.44)
     {$x_0\!\sim\!\mathcal{N}(0,\sigma^2 I)$};
\node[anchor=north, font=\bgfmtiny, align=center, text=oiGreyD]
      at (1.03,7.02) {CTMC prior for\\types, charges};

% -- flow band ----------------------------------------------------------
\shade[left color=oiBlue!7, right color=oiGreen!9, rounded corners=2.2pt]
      (\bandL,\Ay) rectangle (\bandR,\AY);
\draw[rounded corners=2.2pt, draw=oiBlue, line width=0.65pt]
      (\bandL,\Ay) rectangle (\bandR,\AY);

% forward (generative) ODE
\draw[flowarr] (\tz,8.02) -- (\to,8.02);
\node[anchor=south, font=\scriptsize, text=oiBlue!75!black] at (6.32,8.08)
     {$\dfrac{d x_t}{dt}=v_\theta(x_t,t)$\;\; SE(3)-equivariant flow matching};

% the three glyphs along the path
\begin{scope}[shift={(3.40,7.40)}]\bgfmMolNoise\end{scope}
\begin{scope}[shift={(6.32,7.40)}]\bgfmMolMid\end{scope}
\begin{scope}[shift={(9.10,7.40)}]\bgfmMolFinal\end{scope}

% reverse-time (FFJORD) ODE: right -> left, terminates at t = 0
\draw[revarr] (\to,6.98) -- (\tz,6.98);
\node[anchor=north, font=\bgfmtiny, text=oiGreen!50!black] at (6.32,6.94)
     {FFJORD reverse-time ODE $+$ Hutchinson divergence};

% time axis
\draw[oiGrey, line width=0.5pt] (\tz,6.66) -- (\to,6.66);
\foreach \xx in {\tz,\tstar,\to}%
  {\draw[oiGrey, line width=0.5pt] (\xx,6.66) -- (\xx,6.58);}
\node[anchor=north, font=\bgfmtiny] at (\tz,6.57) {$t\!=\!0$};
\node[anchor=north east, font=\bgfmtiny] at (9.60,6.57) {$t^{\star}\!<\!1$};
\node[anchor=north west, font=\bgfmtiny] at (10.13,6.57) {$1$};

% -- x_1 ----------------------------------------------------------------
\draw[bgfmpanel=oiBlue] (10.62,\Ay) rectangle (13.91,\AY);
\node[anchor=north, font=\scriptsize] at (12.26,8.46) {$x_1$};
\node[anchor=north, font=\bgfmtiny, align=center] at (12.26,8.04)
     {positions $+$ atomic numbers\\\textbf{83 elements, bond-free}};
\node[anchor=north, font=\bgfmtiny, align=center, text=oiGreyD]
     at (12.26,7.10) {bonds recovered \emph{post hoc}\\(xyz2mol), never supervised};

% ==================== taps into the two branches =====================
\draw[revarr]  (\tz,\Ay)    -- (\tz,\BY-0.02);
\node[anchor=west, font=\bgfmtiny, text=oiGreen!50!black] at (\tz+0.09,6.20)
     {$\log p_\theta(x_1)$};
\draw[tapearr] (\tstar,\Ay) -- (\tstar,\BY-0.02);
\node[anchor=west, font=\bgfmtiny, text=oiBlue!75!black] at (\tstar+0.09,6.20)
     {$v_\theta(x_{t^\star},t^\star)$};

% ==================== BAND (b): global shape =========================
\draw[bgfmpanel=oiGreen] (0.04,\By) rectangle (6.58,\BY);
\node[ttl, text=oiGreen!45!black] at (0.18,5.76)
     {(b) Global shape: energy consistency};
\node[bdy, anchor=north west] at (0.18,5.58)
     {$\log p_\theta(x_1)=\log p_0(x_0)-\displaystyle\int_{0}^{1}\!\!
       \nabla\!\cdot\!v_\theta\,dt$};

% full-panel-width loss badge: the formula stays on one line
\node[lossbadge=oiGreen, text width=6.06cm, anchor=north west] at (0.18,5.02)
  {$\mathcal{L}_{\mathrm{energy}}=
    \operatorname{Var}_{k}\!\big[\log p_\theta(x_1^{(k)})+E(x_1^{(k)})/kT\big]$\\[1pt]
   {\bgfmtiny the partition function $Z$ cancels exactly}};

% mini potential-energy surface + the K within-parent perturbations
\begin{scope}[shift={(0.30,3.46)}]
  \draw[oiGrey, line width=0.45pt] (0,0) -- (2.26,0);
  \draw[oiGrey, line width=0.45pt] (0,0) -- (0,0.90);
  \node[font=\bgfmtiny, rotate=90, inner sep=1pt] at (-0.14,0.45) {$E$};
  % Boltzmann weight e^{-E/kT} shaded under the curve
  \fill[oiGreen!16]
    plot[smooth] coordinates {(0.12,0.09)(0.44,0.53)(0.76,0.17)(1.08,0.64)
                              (1.42,0.26)(1.76,0.12)(2.10,0.06)}
    -- (2.10,0) -- (0.12,0) -- cycle;
  % E(x): a double well
  \draw[oiBlue, line width=0.8pt]
    plot[smooth] coordinates {(0.12,0.84)(0.44,0.32)(0.76,0.66)(1.08,0.19)
                              (1.42,0.54)(1.76,0.74)(2.10,0.89)};
  \foreach \px/\py in {0.30/0.50, 0.44/0.32, 0.58/0.42, 0.92/0.37,
                       1.08/0.19, 1.24/0.32, 1.54/0.61}
     {\fill[oiGreen!75!black] (\px,\py) circle (1.15pt);}
  \node[anchor=west, font=\bgfmtiny, text=oiGreen!45!black] at (1.14,0.10)
       {$e^{-E/kT}$};
\end{scope}

\node[tny, text width=3.46cm, anchor=north west] at (2.94,4.20)
  {$K$ perturbations of \emph{one} parent molecule\\
   (within-parent variance only)};

% ==================== BAND (c): local gradient =======================
\draw[bgfmpanel=oiVerm] (6.86,\By) rectangle (13.91,\BY);
\node[ttl, text=oiVerm!60!black] at (7.00,5.76) {(c) Local gradient};
\node[bdy, anchor=north west, text width=5.20cm] at (7.00,5.58)
     {closed-form score of the FM marginal\\[3pt]
      $s_\theta(x_t,t)=\dfrac{t\,v_\theta(x_t,t)-x_t}{(1-t)\,\sigma^{2}}$};

% mini inset: the two vectors that L_force aligns
\begin{scope}[shift={(12.30,5.30)}]
  \draw[oiGreyL, line width=0.4pt] (0,0) circle (0.44);
  \fill[oiGrey!85] (0,0) circle (1.7pt);
  \draw[-stealth', line width=0.9pt, draw=oiBlue] (0,0) -- (0.44,0.21);
  \draw[-stealth', line width=0.9pt, draw=oiVerm,
        dash pattern=on 2pt off 1.4pt] (0,0) -- (0.40,0.37);
  \draw[oiGrey, line width=0.35pt] (0.22,0.105) arc (25.6:42.4:0.245);
  \node[anchor=west, font=\bgfmtiny, text=oiBlue!75!black] at (0.46,0.14) {$s_\theta$};
  \node[anchor=west, font=\bgfmtiny, text=oiVerm!70!black] at (0.42,0.44) {$F/kT$};
\end{scope}

\node[lossbadge=oiVerm, text width=6.62cm, anchor=north west] at (7.00,4.60)
  {$\mathcal{L}_{\mathrm{force}}=\big\|\,s_\theta(x_{t^\star}\!,t^\star)
     -F(x_1)/kT\,\big\|_{2}^{2}$\\[1pt]
   {\bgfmtiny at the Boltzmann optimum
    $\nabla\!\log p=-\nabla E/kT=F/kT$;
    \textbf{measured harmful, \Cref{sec:negative}; $\lambda_1\!=\!0$ in every reported model}}};
\node[tny, text width=6.62cm, anchor=north west] at (7.00,3.80)
  {Read out at $t^{\star}\in\{0.85,0.92,0.97\}$ and averaged: the score has a
   $(1-t)^{-1}$ pole at $t\!=\!1$, so the per-atom norm is capped.};

% ==================== BAND (d): the physics bus ======================
\draw[physarr] (3.30,2.96) -- (3.30,3.34);
\node[anchor=west, font=\bgfmtiny, text=oiOrange!45!black] at (3.38,3.16)
     {$E(x_1^{(k)})/kT$};
\draw[physarr] (10.20,2.96) -- (10.20,3.34);
\node[anchor=west, font=\bgfmtiny, text=oiOrange!45!black] at (10.28,3.16)
     {$F(x_1)/kT=-\nabla E/kT$};

\draw[rounded corners=2.2pt, draw=oiOrange!85!black, line width=0.75pt,
      fill=oiOrange!12] (0.04,2.10) rectangle (13.91,2.92);
\node[anchor=west, font=\scriptsize\bfseries, text=oiOrange!35!black]
     at (0.18,2.66)
     {(d) Universal neural potential $E(x)$ --- OMol25 / eSEN};
\node[anchor=west, font=\bgfmtiny, text=oiGreyD] at (0.18,2.31)
     {3.9M molecules, 205M atoms, 83 elements; DFT forces \emph{and}
      energies; \textbf{no bond labels} --- one oracle, both terms};

% ==================== BAND (e): total loss + legend ==================
\draw[rounded corners=2pt, draw=oiGrey, line width=0.55pt, fill=white]
      (2.30,1.06) rectangle (11.66,1.82);
\node[anchor=center, font=\small] at (6.98,1.44)
  {$\mathcal{L}_{\mathrm{BGFM}}=\mathcal{L}_{\mathrm{FM}}
    +\lambda_1\,\textcolor{oiVerm}{\mathcal{L}_{\mathrm{force}}}
    +\lambda_2\,\textcolor{oiGreen}{\mathcal{L}_{\mathrm{energy}}}
    \qquad
    \textcolor{oiGreen}{\mathcal{L}_{\mathrm{energy}}}\!=\!0
    \iff p_\theta(\cdot\mid c)=\pi(\cdot\mid c)$};

\begin{scope}[shift={(0.30,0.28)}]
  \def\lw{0.26}
  \foreach \i/\c/\lbl in {%
      0/oiBlue/{model, flow matching},%
      1/oiGreen/{$\mathcal{L}_{\mathrm{energy}}$: global shape},%
      2/oiVerm/{$\mathcal{L}_{\mathrm{force}}$: local gradient},%
      3/oiOrange/{physics oracle (DFT labels)}}%
  {\pgfmathsetmacro{\xo}{\i*3.30}%
   \fill[\c!22] (\xo,0) rectangle (\xo+\lw,0.24);%
   \draw[\c, line width=0.6pt] (\xo,0) rectangle (\xo+\lw,0.24);%
   \node[anchor=west, font=\bgfmtiny] at (\xo+\lw+0.10,0.12) {\lbl};}
\end{scope}

\end{tikzpicture}

\bgfmFinish

%           \caption{...}\label{fig:method}
%         \end{figure}
%
%  SIZE.  Natural size 13.95cm x 8.70cm == 5.49in x 3.43in.
%  ICLR is a SINGLE 5.5in-wide text column (10pt Times, 11pt leading), so
%  this is a full-text-width figure at scale 1.0.  DO NOT rescale it: every
%  glyph inside is \footnotesize / \scriptsize / \tiny of the 10pt body
%  font, and scaling desynchronises in-figure text from the caption.
%  If a shorter figure is needed, delete band (e) (the total-loss box and
%  legend, y < 1.7) and state the loss in the caption instead -- that saves
%  1.5cm without touching any other coordinate.
%
%  LAYOUT (five bands, top to bottom).  Wire routing is deliberately
%  crossing-free:
%    (a) generative spine   prior noise -> flow matching -> x_1
%    (b) global-shape term  FFJORD log p  vs  E/kT     -> L_energy  [LEFT]
%    (c) local-gradient term closed-form score vs F/kT -> L_force   [RIGHT]
%    (d) physics bus        universal neural potential (OMol25 / eSEN)
%    (e) total loss + legend
%  (b) is on the left because the reverse-time (FFJORD) arrow terminates at
%  t=0 on the left; (c) is on the right because the score is read out at
%  t*=0.95 on the right; (d) is a full-width bus underneath so both terms
%  are fed by short vertical arrows and no wire ever crosses another.
% =====================================================================

% ---- standalone bootstrap (skipped when \input into a tikz-enabled doc) ----
\ifx\tikzpicture\undefined
  \documentclass[10pt]{article}
  \usepackage[paperwidth=14.30cm,paperheight=9.15cm,margin=0.15cm,noheadfoot]{geometry}
  \usepackage{amsmath,amssymb,bm}
  \usepackage{tikz}
  \setlength{\topskip}{0pt}
  \setlength{\parindent}{0pt}
  \pagestyle{empty}
  \begin{document}%
  \def\bgfmFinish{\end{document}}
\fi
\providecommand{\bgfmFinish}{\relax}
% \pending{...}: red placeholder marker, shared with the main text so that a
% single grep finds every not-yet-real number in the manuscript.
\providecommand{\pending}[1]{\textcolor{red}{#1}}
\providecommand{\bgfmFigScale}{1.0}

% 6pt/6.6pt: one step above \tiny (5pt at 10pt base), which is below the
% conventional 6pt floor for in-figure text.
\providecommand{\bgfmtiny}{\fontsize{6}{6.6}\selectfont}
\usetikzlibrary{positioning,calc,fit,backgrounds,decorations.pathmorphing,%
                arrows,shapes.geometric}

% ---------------------------------------------------------------- palette
% Okabe & Ito (2008): the de-facto standard CVD-safe qualitative set.
\definecolor{oiOrange}{HTML}{E69F00}
\definecolor{oiSky}   {HTML}{56B4E9}
\definecolor{oiGreen} {HTML}{009E73}
\definecolor{oiYellow}{HTML}{F0E442}
\definecolor{oiBlue}  {HTML}{0072B2}
\definecolor{oiVerm}  {HTML}{D55E00}
\definecolor{oiPurple}{HTML}{CC79A7}
\definecolor{oiGrey}  {HTML}{7F7F7F}
\definecolor{oiGreyL} {HTML}{BFBFBF}
\definecolor{oiGreyD} {HTML}{555555}
%
% SEMANTIC assignment -- identical in this figure and in the matplotlib
% figures (make_experiment_figures.py::PALETTE), so the paper reads as one
% visual system:
%   oiBlue    model / flow matching / learned quantities
%   oiGreen   L_energy  (global shape, FFJORD, Boltzmann)
%   oiVerm    L_force   (local gradient, score--force matching)
%   oiOrange  physics oracle (universal neural potential, DFT labels)
%   oiPurple  negative control (Y-scrambled)
%   oiGrey    no-physics baseline / neutral structure

% ---------------------------------------------------------------- styles
\tikzset{
  bgfmpanel/.style={rounded corners=2.2pt, draw=#1, line width=0.65pt,
                    fill=#1!5},
  ttl/.style={anchor=west, font=\footnotesize\bfseries, inner sep=0pt},
  bdy/.style={anchor=west, font=\scriptsize, inner sep=0pt, align=left},
  tny/.style={anchor=west, font=\bgfmtiny, inner sep=0pt, align=left,
              text=oiGreyD},
  flowarr/.style={-stealth', line width=0.9pt, draw=oiBlue},
  revarr/.style={-stealth', line width=0.75pt, draw=oiGreen,
                 dash pattern=on 2.2pt off 1.6pt},
  physarr/.style={-stealth', line width=0.8pt, draw=oiOrange!85!black},
  tapearr/.style={-stealth', line width=0.8pt, draw=oiBlue},
  lossbadge/.style={rounded corners=2pt, draw=#1, line width=0.6pt,
                    fill=#1!12, inner sep=2.6pt, align=left,
                    font=\scriptsize},
}

% ------------------------------------------------------- molecule glyphs
% Deterministic coordinates (no \pgfmathrandom) so the figure is
% byte-reproducible across compiles and across TeX distributions.
\newcommand{\bgfmAtom}[3]{\fill[#3] (#1,#2) circle (1.45pt);%
  \draw[oiGrey!60, line width=0.15pt] (#1,#2) circle (1.45pt);}
\newcommand{\bgfmAtomS}[3]{\fill[#3] (#1,#2) circle (0.95pt);%
  \draw[oiGrey!60, line width=0.15pt] (#1,#2) circle (0.95pt);}
\newcommand{\bgfmBond}[4]{\draw[oiGrey!85, line width=0.55pt] (#1,#2) -- (#3,#4);}

% t ~ 0 : unstructured point cloud, no bonds
\newcommand{\bgfmMolNoise}{%
  \bgfmAtomS{-0.34}{ 0.30}{oiGreyL}\bgfmAtomS{ 0.10}{ 0.36}{oiGreyL}%
  \bgfmAtomS{ 0.36}{ 0.06}{oiGreyL}\bgfmAtomS{-0.10}{-0.02}{oiGreyL}%
  \bgfmAtomS{-0.40}{-0.26}{oiGreyL}\bgfmAtomS{ 0.06}{-0.34}{oiGreyL}%
  \bgfmAtomS{ 0.40}{-0.32}{oiGreyL}\bgfmAtomS{ 0.24}{ 0.22}{oiGreyL}%
  \bgfmAtomS{-0.18}{ 0.12}{oiGreyL}}

% intermediate t : partially condensed, proto-structure
\newcommand{\bgfmMolMid}{%
  \bgfmBond{-0.20}{0.14}{0.04}{0.22}\bgfmBond{0.04}{0.22}{0.24}{0.02}%
  \bgfmBond{-0.20}{0.14}{-0.14}{-0.16}%
  \bgfmAtomS{-0.20}{ 0.14}{oiGrey!70}\bgfmAtomS{ 0.04}{ 0.22}{oiGrey!70}%
  \bgfmAtomS{ 0.24}{ 0.02}{oiGrey!70}\bgfmAtomS{-0.14}{-0.16}{oiGrey!70}%
  \bgfmAtomS{ 0.30}{ 0.30}{oiGreyL}\bgfmAtomS{-0.36}{-0.24}{oiGreyL}%
  \bgfmAtomS{ 0.20}{-0.30}{oiGreyL}}

% t = 1 : a clean structure -- 5-ring + N + O + a transition metal (green),
% so the "83 elements / transition metals" claim is visible in the glyph.
\newcommand{\bgfmMolFinal}{%
  \bgfmBond{0.000}{0.280}{-0.266}{0.087}\bgfmBond{-0.266}{0.087}{-0.165}{-0.227}%
  \bgfmBond{-0.165}{-0.227}{0.165}{-0.227}\bgfmBond{0.165}{-0.227}{0.266}{0.087}%
  \bgfmBond{0.266}{0.087}{0.000}{0.280}%
  \bgfmBond{0.266}{0.087}{0.560}{0.250}%
  \bgfmBond{-0.266}{0.087}{-0.520}{0.300}%
  \bgfmBond{-0.165}{-0.227}{-0.360}{-0.440}%
  \bgfmAtom{ 0.000}{ 0.280}{oiGrey!85}\bgfmAtom{-0.266}{ 0.087}{oiGrey!85}%
  \bgfmAtom{-0.165}{-0.227}{oiGrey!85}\bgfmAtom{ 0.165}{-0.227}{oiGrey!85}%
  \bgfmAtom{ 0.266}{ 0.087}{oiBlue}%
  \fill[oiGreen] (0.560,0.250) circle (2.4pt);%
  \draw[oiGrey!70, line width=0.2pt] (0.560,0.250) circle (2.4pt);%
  \bgfmAtomS{-0.520}{ 0.300}{oiVerm}%
  \bgfmAtomS{-0.360}{-0.440}{white}}

% =====================================================================
\begin{tikzpicture}[x=1cm, y=1cm, scale=\bgfmFigScale, >=stealth']
% guarantee a stable bounding box
\path (0,0.06) rectangle (13.95,8.66);

% ---- shared geometry -------------------------------------------------
\def\Ay{6.42}\def\AY{8.52}          % band (a) vertical extent
\def\bandL{2.30}\def\bandR{10.34}   % flow band horizontal extent
\def\tz{2.55}\def\to{10.09}         % x of t=0 and t=1 on the time axis
\def\tstar{9.71}                    % x of t* = 0.95
\def\By{3.36}\def\BY{5.96}          % bands (b)/(c) vertical extent

% ==================== BAND (a): generative spine =====================
% -- prior --------------------------------------------------------------
\draw[bgfmpanel=oiBlue] (0.04,\Ay) rectangle (2.02,\AY);
\node[anchor=north west, font=\scriptsize\bfseries] at (0.12,8.48) {(a)};
\begin{scope}[shift={(1.06,7.88)}]\bgfmMolNoise\end{scope}
\node[anchor=north, font=\scriptsize] at (1.03,7.44)
     {$x_0\!\sim\!\mathcal{N}(0,\sigma^2 I)$};
\node[anchor=north, font=\bgfmtiny, align=center, text=oiGreyD]
      at (1.03,7.02) {CTMC prior for\\types, charges};

% -- flow band ----------------------------------------------------------
\shade[left color=oiBlue!7, right color=oiGreen!9, rounded corners=2.2pt]
      (\bandL,\Ay) rectangle (\bandR,\AY);
\draw[rounded corners=2.2pt, draw=oiBlue, line width=0.65pt]
      (\bandL,\Ay) rectangle (\bandR,\AY);

% forward (generative) ODE
\draw[flowarr] (\tz,8.02) -- (\to,8.02);
\node[anchor=south, font=\scriptsize, text=oiBlue!75!black] at (6.32,8.08)
     {$\dfrac{d x_t}{dt}=v_\theta(x_t,t)$\;\; SE(3)-equivariant flow matching};

% the three glyphs along the path
\begin{scope}[shift={(3.40,7.40)}]\bgfmMolNoise\end{scope}
\begin{scope}[shift={(6.32,7.40)}]\bgfmMolMid\end{scope}
\begin{scope}[shift={(9.10,7.40)}]\bgfmMolFinal\end{scope}

% reverse-time (FFJORD) ODE: right -> left, terminates at t = 0
\draw[revarr] (\to,6.98) -- (\tz,6.98);
\node[anchor=north, font=\bgfmtiny, text=oiGreen!50!black] at (6.32,6.94)
     {FFJORD reverse-time ODE $+$ Hutchinson divergence};

% time axis
\draw[oiGrey, line width=0.5pt] (\tz,6.66) -- (\to,6.66);
\foreach \xx in {\tz,\tstar,\to}%
  {\draw[oiGrey, line width=0.5pt] (\xx,6.66) -- (\xx,6.58);}
\node[anchor=north, font=\bgfmtiny] at (\tz,6.57) {$t\!=\!0$};
\node[anchor=north east, font=\bgfmtiny] at (9.60,6.57) {$t^{\star}\!<\!1$};
\node[anchor=north west, font=\bgfmtiny] at (10.13,6.57) {$1$};

% -- x_1 ----------------------------------------------------------------
\draw[bgfmpanel=oiBlue] (10.62,\Ay) rectangle (13.91,\AY);
\node[anchor=north, font=\scriptsize] at (12.26,8.46) {$x_1$};
\node[anchor=north, font=\bgfmtiny, align=center] at (12.26,8.04)
     {positions $+$ atomic numbers\\\textbf{83 elements, bond-free}};
\node[anchor=north, font=\bgfmtiny, align=center, text=oiGreyD]
     at (12.26,7.10) {bonds recovered \emph{post hoc}\\(xyz2mol), never supervised};

% ==================== taps into the two branches =====================
\draw[revarr]  (\tz,\Ay)    -- (\tz,\BY-0.02);
\node[anchor=west, font=\bgfmtiny, text=oiGreen!50!black] at (\tz+0.09,6.20)
     {$\log p_\theta(x_1)$};
\draw[tapearr] (\tstar,\Ay) -- (\tstar,\BY-0.02);
\node[anchor=west, font=\bgfmtiny, text=oiBlue!75!black] at (\tstar+0.09,6.20)
     {$v_\theta(x_{t^\star},t^\star)$};

% ==================== BAND (b): global shape =========================
\draw[bgfmpanel=oiGreen] (0.04,\By) rectangle (6.58,\BY);
\node[ttl, text=oiGreen!45!black] at (0.18,5.76)
     {(b) Global shape: energy consistency};
\node[bdy, anchor=north west] at (0.18,5.58)
     {$\log p_\theta(x_1)=\log p_0(x_0)-\displaystyle\int_{0}^{1}\!\!
       \nabla\!\cdot\!v_\theta\,dt$};

% full-panel-width loss badge: the formula stays on one line
\node[lossbadge=oiGreen, text width=6.06cm, anchor=north west] at (0.18,5.02)
  {$\mathcal{L}_{\mathrm{energy}}=
    \operatorname{Var}_{k}\!\big[\log p_\theta(x_1^{(k)})+E(x_1^{(k)})/kT\big]$\\[1pt]
   {\bgfmtiny the partition function $Z$ cancels exactly}};

% mini potential-energy surface + the K within-parent perturbations
\begin{scope}[shift={(0.30,3.46)}]
  \draw[oiGrey, line width=0.45pt] (0,0) -- (2.26,0);
  \draw[oiGrey, line width=0.45pt] (0,0) -- (0,0.90);
  \node[font=\bgfmtiny, rotate=90, inner sep=1pt] at (-0.14,0.45) {$E$};
  % Boltzmann weight e^{-E/kT} shaded under the curve
  \fill[oiGreen!16]
    plot[smooth] coordinates {(0.12,0.09)(0.44,0.53)(0.76,0.17)(1.08,0.64)
                              (1.42,0.26)(1.76,0.12)(2.10,0.06)}
    -- (2.10,0) -- (0.12,0) -- cycle;
  % E(x): a double well
  \draw[oiBlue, line width=0.8pt]
    plot[smooth] coordinates {(0.12,0.84)(0.44,0.32)(0.76,0.66)(1.08,0.19)
                              (1.42,0.54)(1.76,0.74)(2.10,0.89)};
  \foreach \px/\py in {0.30/0.50, 0.44/0.32, 0.58/0.42, 0.92/0.37,
                       1.08/0.19, 1.24/0.32, 1.54/0.61}
     {\fill[oiGreen!75!black] (\px,\py) circle (1.15pt);}
  \node[anchor=west, font=\bgfmtiny, text=oiGreen!45!black] at (1.14,0.10)
       {$e^{-E/kT}$};
\end{scope}

\node[tny, text width=3.46cm, anchor=north west] at (2.94,4.20)
  {$K$ perturbations of \emph{one} parent molecule\\
   (within-parent variance only)};

% ==================== BAND (c): local gradient =======================
\draw[bgfmpanel=oiVerm] (6.86,\By) rectangle (13.91,\BY);
\node[ttl, text=oiVerm!60!black] at (7.00,5.76) {(c) Local gradient};
\node[bdy, anchor=north west, text width=5.20cm] at (7.00,5.58)
     {closed-form score of the FM marginal\\[3pt]
      $s_\theta(x_t,t)=\dfrac{t\,v_\theta(x_t,t)-x_t}{(1-t)\,\sigma^{2}}$};

% mini inset: the two vectors that L_force aligns
\begin{scope}[shift={(12.30,5.30)}]
  \draw[oiGreyL, line width=0.4pt] (0,0) circle (0.44);
  \fill[oiGrey!85] (0,0) circle (1.7pt);
  \draw[-stealth', line width=0.9pt, draw=oiBlue] (0,0) -- (0.44,0.21);
  \draw[-stealth', line width=0.9pt, draw=oiVerm,
        dash pattern=on 2pt off 1.4pt] (0,0) -- (0.40,0.37);
  \draw[oiGrey, line width=0.35pt] (0.22,0.105) arc (25.6:42.4:0.245);
  \node[anchor=west, font=\bgfmtiny, text=oiBlue!75!black] at (0.46,0.14) {$s_\theta$};
  \node[anchor=west, font=\bgfmtiny, text=oiVerm!70!black] at (0.42,0.44) {$F/kT$};
\end{scope}

\node[lossbadge=oiVerm, text width=6.62cm, anchor=north west] at (7.00,4.60)
  {$\mathcal{L}_{\mathrm{force}}=\big\|\,s_\theta(x_{t^\star}\!,t^\star)
     -F(x_1)/kT\,\big\|_{2}^{2}$\\[1pt]
   {\bgfmtiny at the Boltzmann optimum
    $\nabla\!\log p=-\nabla E/kT=F/kT$;
    \textbf{measured harmful, \Cref{sec:negative}; $\lambda_1\!=\!0$ in every reported model}}};
\node[tny, text width=6.62cm, anchor=north west] at (7.00,3.80)
  {Read out at $t^{\star}\in\{0.85,0.92,0.97\}$ and averaged: the score has a
   $(1-t)^{-1}$ pole at $t\!=\!1$, so the per-atom norm is capped.};

% ==================== BAND (d): the physics bus ======================
\draw[physarr] (3.30,2.96) -- (3.30,3.34);
\node[anchor=west, font=\bgfmtiny, text=oiOrange!45!black] at (3.38,3.16)
     {$E(x_1^{(k)})/kT$};
\draw[physarr] (10.20,2.96) -- (10.20,3.34);
\node[anchor=west, font=\bgfmtiny, text=oiOrange!45!black] at (10.28,3.16)
     {$F(x_1)/kT=-\nabla E/kT$};

\draw[rounded corners=2.2pt, draw=oiOrange!85!black, line width=0.75pt,
      fill=oiOrange!12] (0.04,2.10) rectangle (13.91,2.92);
\node[anchor=west, font=\scriptsize\bfseries, text=oiOrange!35!black]
     at (0.18,2.66)
     {(d) Universal neural potential $E(x)$ --- OMol25 / eSEN};
\node[anchor=west, font=\bgfmtiny, text=oiGreyD] at (0.18,2.31)
     {3.9M molecules, 205M atoms, 83 elements; DFT forces \emph{and}
      energies; \textbf{no bond labels} --- one oracle, both terms};

% ==================== BAND (e): total loss + legend ==================
\draw[rounded corners=2pt, draw=oiGrey, line width=0.55pt, fill=white]
      (2.30,1.06) rectangle (11.66,1.82);
\node[anchor=center, font=\small] at (6.98,1.44)
  {$\mathcal{L}_{\mathrm{BGFM}}=\mathcal{L}_{\mathrm{FM}}
    +\lambda_1\,\textcolor{oiVerm}{\mathcal{L}_{\mathrm{force}}}
    +\lambda_2\,\textcolor{oiGreen}{\mathcal{L}_{\mathrm{energy}}}
    \qquad
    \textcolor{oiGreen}{\mathcal{L}_{\mathrm{energy}}}\!=\!0
    \iff p_\theta(\cdot\mid c)=\pi(\cdot\mid c)$};

\begin{scope}[shift={(0.30,0.28)}]
  \def\lw{0.26}
  \foreach \i/\c/\lbl in {%
      0/oiBlue/{model, flow matching},%
      1/oiGreen/{$\mathcal{L}_{\mathrm{energy}}$: global shape},%
      2/oiVerm/{$\mathcal{L}_{\mathrm{force}}$: local gradient},%
      3/oiOrange/{physics oracle (DFT labels)}}%
  {\pgfmathsetmacro{\xo}{\i*3.30}%
   \fill[\c!22] (\xo,0) rectangle (\xo+\lw,0.24);%
   \draw[\c, line width=0.6pt] (\xo,0) rectangle (\xo+\lw,0.24);%
   \node[anchor=west, font=\bgfmtiny] at (\xo+\lw+0.10,0.12) {\lbl};}
\end{scope}

\end{tikzpicture}

\bgfmFinish

%           \caption{...}\label{fig:method}
%         \end{figure}
%
%  SIZE.  Natural size 13.95cm x 8.70cm == 5.49in x 3.43in.
%  ICLR is a SINGLE 5.5in-wide text column (10pt Times, 11pt leading), so
%  this is a full-text-width figure at scale 1.0.  DO NOT rescale it: every
%  glyph inside is \footnotesize / \scriptsize / \tiny of the 10pt body
%  font, and scaling desynchronises in-figure text from the caption.
%  If a shorter figure is needed, delete band (e) (the total-loss box and
%  legend, y < 1.7) and state the loss in the caption instead -- that saves
%  1.5cm without touching any other coordinate.
%
%  LAYOUT (five bands, top to bottom).  Wire routing is deliberately
%  crossing-free:
%    (a) generative spine   prior noise -> flow matching -> x_1
%    (b) global-shape term  FFJORD log p  vs  E/kT     -> L_energy  [LEFT]
%    (c) local-gradient term closed-form score vs F/kT -> L_force   [RIGHT]
%    (d) physics bus        universal neural potential (OMol25 / eSEN)
%    (e) total loss + legend
%  (b) is on the left because the reverse-time (FFJORD) arrow terminates at
%  t=0 on the left; (c) is on the right because the score is read out at
%  t*=0.95 on the right; (d) is a full-width bus underneath so both terms
%  are fed by short vertical arrows and no wire ever crosses another.
% =====================================================================

% ---- standalone bootstrap (skipped when \input into a tikz-enabled doc) ----
\ifx\tikzpicture\undefined
  \documentclass[10pt]{article}
  \usepackage[paperwidth=14.30cm,paperheight=9.15cm,margin=0.15cm,noheadfoot]{geometry}
  \usepackage{amsmath,amssymb,bm}
  \usepackage{tikz}
  \setlength{\topskip}{0pt}
  \setlength{\parindent}{0pt}
  \pagestyle{empty}
  \begin{document}%
  \def\bgfmFinish{\end{document}}
\fi
\providecommand{\bgfmFinish}{\relax}
% \pending{...}: red placeholder marker, shared with the main text so that a
% single grep finds every not-yet-real number in the manuscript.
\providecommand{\pending}[1]{\textcolor{red}{#1}}
\providecommand{\bgfmFigScale}{1.0}

% 6pt/6.6pt: one step above \tiny (5pt at 10pt base), which is below the
% conventional 6pt floor for in-figure text.
\providecommand{\bgfmtiny}{\fontsize{6}{6.6}\selectfont}
\usetikzlibrary{positioning,calc,fit,backgrounds,decorations.pathmorphing,%
                arrows,shapes.geometric}

% ---------------------------------------------------------------- palette
% Okabe & Ito (2008): the de-facto standard CVD-safe qualitative set.
\definecolor{oiOrange}{HTML}{E69F00}
\definecolor{oiSky}   {HTML}{56B4E9}
\definecolor{oiGreen} {HTML}{009E73}
\definecolor{oiYellow}{HTML}{F0E442}
\definecolor{oiBlue}  {HTML}{0072B2}
\definecolor{oiVerm}  {HTML}{D55E00}
\definecolor{oiPurple}{HTML}{CC79A7}
\definecolor{oiGrey}  {HTML}{7F7F7F}
\definecolor{oiGreyL} {HTML}{BFBFBF}
\definecolor{oiGreyD} {HTML}{555555}
%
% SEMANTIC assignment -- identical in this figure and in the matplotlib
% figures (make_experiment_figures.py::PALETTE), so the paper reads as one
% visual system:
%   oiBlue    model / flow matching / learned quantities
%   oiGreen   L_energy  (global shape, FFJORD, Boltzmann)
%   oiVerm    L_force   (local gradient, score--force matching)
%   oiOrange  physics oracle (universal neural potential, DFT labels)
%   oiPurple  negative control (Y-scrambled)
%   oiGrey    no-physics baseline / neutral structure

% ---------------------------------------------------------------- styles
\tikzset{
  bgfmpanel/.style={rounded corners=2.2pt, draw=#1, line width=0.65pt,
                    fill=#1!5},
  ttl/.style={anchor=west, font=\footnotesize\bfseries, inner sep=0pt},
  bdy/.style={anchor=west, font=\scriptsize, inner sep=0pt, align=left},
  tny/.style={anchor=west, font=\bgfmtiny, inner sep=0pt, align=left,
              text=oiGreyD},
  flowarr/.style={-stealth', line width=0.9pt, draw=oiBlue},
  revarr/.style={-stealth', line width=0.75pt, draw=oiGreen,
                 dash pattern=on 2.2pt off 1.6pt},
  physarr/.style={-stealth', line width=0.8pt, draw=oiOrange!85!black},
  tapearr/.style={-stealth', line width=0.8pt, draw=oiBlue},
  lossbadge/.style={rounded corners=2pt, draw=#1, line width=0.6pt,
                    fill=#1!12, inner sep=2.6pt, align=left,
                    font=\scriptsize},
}

% ------------------------------------------------------- molecule glyphs
% Deterministic coordinates (no \pgfmathrandom) so the figure is
% byte-reproducible across compiles and across TeX distributions.
\newcommand{\bgfmAtom}[3]{\fill[#3] (#1,#2) circle (1.45pt);%
  \draw[oiGrey!60, line width=0.15pt] (#1,#2) circle (1.45pt);}
\newcommand{\bgfmAtomS}[3]{\fill[#3] (#1,#2) circle (0.95pt);%
  \draw[oiGrey!60, line width=0.15pt] (#1,#2) circle (0.95pt);}
\newcommand{\bgfmBond}[4]{\draw[oiGrey!85, line width=0.55pt] (#1,#2) -- (#3,#4);}

% t ~ 0 : unstructured point cloud, no bonds
\newcommand{\bgfmMolNoise}{%
  \bgfmAtomS{-0.34}{ 0.30}{oiGreyL}\bgfmAtomS{ 0.10}{ 0.36}{oiGreyL}%
  \bgfmAtomS{ 0.36}{ 0.06}{oiGreyL}\bgfmAtomS{-0.10}{-0.02}{oiGreyL}%
  \bgfmAtomS{-0.40}{-0.26}{oiGreyL}\bgfmAtomS{ 0.06}{-0.34}{oiGreyL}%
  \bgfmAtomS{ 0.40}{-0.32}{oiGreyL}\bgfmAtomS{ 0.24}{ 0.22}{oiGreyL}%
  \bgfmAtomS{-0.18}{ 0.12}{oiGreyL}}

% intermediate t : partially condensed, proto-structure
\newcommand{\bgfmMolMid}{%
  \bgfmBond{-0.20}{0.14}{0.04}{0.22}\bgfmBond{0.04}{0.22}{0.24}{0.02}%
  \bgfmBond{-0.20}{0.14}{-0.14}{-0.16}%
  \bgfmAtomS{-0.20}{ 0.14}{oiGrey!70}\bgfmAtomS{ 0.04}{ 0.22}{oiGrey!70}%
  \bgfmAtomS{ 0.24}{ 0.02}{oiGrey!70}\bgfmAtomS{-0.14}{-0.16}{oiGrey!70}%
  \bgfmAtomS{ 0.30}{ 0.30}{oiGreyL}\bgfmAtomS{-0.36}{-0.24}{oiGreyL}%
  \bgfmAtomS{ 0.20}{-0.30}{oiGreyL}}

% t = 1 : a clean structure -- 5-ring + N + O + a transition metal (green),
% so the "83 elements / transition metals" claim is visible in the glyph.
\newcommand{\bgfmMolFinal}{%
  \bgfmBond{0.000}{0.280}{-0.266}{0.087}\bgfmBond{-0.266}{0.087}{-0.165}{-0.227}%
  \bgfmBond{-0.165}{-0.227}{0.165}{-0.227}\bgfmBond{0.165}{-0.227}{0.266}{0.087}%
  \bgfmBond{0.266}{0.087}{0.000}{0.280}%
  \bgfmBond{0.266}{0.087}{0.560}{0.250}%
  \bgfmBond{-0.266}{0.087}{-0.520}{0.300}%
  \bgfmBond{-0.165}{-0.227}{-0.360}{-0.440}%
  \bgfmAtom{ 0.000}{ 0.280}{oiGrey!85}\bgfmAtom{-0.266}{ 0.087}{oiGrey!85}%
  \bgfmAtom{-0.165}{-0.227}{oiGrey!85}\bgfmAtom{ 0.165}{-0.227}{oiGrey!85}%
  \bgfmAtom{ 0.266}{ 0.087}{oiBlue}%
  \fill[oiGreen] (0.560,0.250) circle (2.4pt);%
  \draw[oiGrey!70, line width=0.2pt] (0.560,0.250) circle (2.4pt);%
  \bgfmAtomS{-0.520}{ 0.300}{oiVerm}%
  \bgfmAtomS{-0.360}{-0.440}{white}}

% =====================================================================
\begin{tikzpicture}[x=1cm, y=1cm, scale=\bgfmFigScale, >=stealth']
% guarantee a stable bounding box
\path (0,0.06) rectangle (13.95,8.66);

% ---- shared geometry -------------------------------------------------
\def\Ay{6.42}\def\AY{8.52}          % band (a) vertical extent
\def\bandL{2.30}\def\bandR{10.34}   % flow band horizontal extent
\def\tz{2.55}\def\to{10.09}         % x of t=0 and t=1 on the time axis
\def\tstar{9.71}                    % x of t* = 0.95
\def\By{3.36}\def\BY{5.96}          % bands (b)/(c) vertical extent

% ==================== BAND (a): generative spine =====================
% -- prior --------------------------------------------------------------
\draw[bgfmpanel=oiBlue] (0.04,\Ay) rectangle (2.02,\AY);
\node[anchor=north west, font=\scriptsize\bfseries] at (0.12,8.48) {(a)};
\begin{scope}[shift={(1.06,7.88)}]\bgfmMolNoise\end{scope}
\node[anchor=north, font=\scriptsize] at (1.03,7.44)
     {$x_0\!\sim\!\mathcal{N}(0,\sigma^2 I)$};
\node[anchor=north, font=\bgfmtiny, align=center, text=oiGreyD]
      at (1.03,7.02) {CTMC prior for\\types, charges};

% -- flow band ----------------------------------------------------------
\shade[left color=oiBlue!7, right color=oiGreen!9, rounded corners=2.2pt]
      (\bandL,\Ay) rectangle (\bandR,\AY);
\draw[rounded corners=2.2pt, draw=oiBlue, line width=0.65pt]
      (\bandL,\Ay) rectangle (\bandR,\AY);

% forward (generative) ODE
\draw[flowarr] (\tz,8.02) -- (\to,8.02);
\node[anchor=south, font=\scriptsize, text=oiBlue!75!black] at (6.32,8.08)
     {$\dfrac{d x_t}{dt}=v_\theta(x_t,t)$\;\; SE(3)-equivariant flow matching};

% the three glyphs along the path
\begin{scope}[shift={(3.40,7.40)}]\bgfmMolNoise\end{scope}
\begin{scope}[shift={(6.32,7.40)}]\bgfmMolMid\end{scope}
\begin{scope}[shift={(9.10,7.40)}]\bgfmMolFinal\end{scope}

% reverse-time (FFJORD) ODE: right -> left, terminates at t = 0
\draw[revarr] (\to,6.98) -- (\tz,6.98);
\node[anchor=north, font=\bgfmtiny, text=oiGreen!50!black] at (6.32,6.94)
     {FFJORD reverse-time ODE $+$ Hutchinson divergence};

% time axis
\draw[oiGrey, line width=0.5pt] (\tz,6.66) -- (\to,6.66);
\foreach \xx in {\tz,\tstar,\to}%
  {\draw[oiGrey, line width=0.5pt] (\xx,6.66) -- (\xx,6.58);}
\node[anchor=north, font=\bgfmtiny] at (\tz,6.57) {$t\!=\!0$};
\node[anchor=north east, font=\bgfmtiny] at (9.60,6.57) {$t^{\star}\!<\!1$};
\node[anchor=north west, font=\bgfmtiny] at (10.13,6.57) {$1$};

% -- x_1 ----------------------------------------------------------------
\draw[bgfmpanel=oiBlue] (10.62,\Ay) rectangle (13.91,\AY);
\node[anchor=north, font=\scriptsize] at (12.26,8.46) {$x_1$};
\node[anchor=north, font=\bgfmtiny, align=center] at (12.26,8.04)
     {positions $+$ atomic numbers\\\textbf{83 elements, bond-free}};
\node[anchor=north, font=\bgfmtiny, align=center, text=oiGreyD]
     at (12.26,7.10) {bonds recovered \emph{post hoc}\\(xyz2mol), never supervised};

% ==================== taps into the two branches =====================
\draw[revarr]  (\tz,\Ay)    -- (\tz,\BY-0.02);
\node[anchor=west, font=\bgfmtiny, text=oiGreen!50!black] at (\tz+0.09,6.20)
     {$\log p_\theta(x_1)$};
\draw[tapearr] (\tstar,\Ay) -- (\tstar,\BY-0.02);
\node[anchor=west, font=\bgfmtiny, text=oiBlue!75!black] at (\tstar+0.09,6.20)
     {$v_\theta(x_{t^\star},t^\star)$};

% ==================== BAND (b): global shape =========================
\draw[bgfmpanel=oiGreen] (0.04,\By) rectangle (6.58,\BY);
\node[ttl, text=oiGreen!45!black] at (0.18,5.76)
     {(b) Global shape: energy consistency};
\node[bdy, anchor=north west] at (0.18,5.58)
     {$\log p_\theta(x_1)=\log p_0(x_0)-\displaystyle\int_{0}^{1}\!\!
       \nabla\!\cdot\!v_\theta\,dt$};

% full-panel-width loss badge: the formula stays on one line
\node[lossbadge=oiGreen, text width=6.06cm, anchor=north west] at (0.18,5.02)
  {$\mathcal{L}_{\mathrm{energy}}=
    \operatorname{Var}_{k}\!\big[\log p_\theta(x_1^{(k)})+E(x_1^{(k)})/kT\big]$\\[1pt]
   {\bgfmtiny the partition function $Z$ cancels exactly}};

% mini potential-energy surface + the K within-parent perturbations
\begin{scope}[shift={(0.30,3.46)}]
  \draw[oiGrey, line width=0.45pt] (0,0) -- (2.26,0);
  \draw[oiGrey, line width=0.45pt] (0,0) -- (0,0.90);
  \node[font=\bgfmtiny, rotate=90, inner sep=1pt] at (-0.14,0.45) {$E$};
  % Boltzmann weight e^{-E/kT} shaded under the curve
  \fill[oiGreen!16]
    plot[smooth] coordinates {(0.12,0.09)(0.44,0.53)(0.76,0.17)(1.08,0.64)
                              (1.42,0.26)(1.76,0.12)(2.10,0.06)}
    -- (2.10,0) -- (0.12,0) -- cycle;
  % E(x): a double well
  \draw[oiBlue, line width=0.8pt]
    plot[smooth] coordinates {(0.12,0.84)(0.44,0.32)(0.76,0.66)(1.08,0.19)
                              (1.42,0.54)(1.76,0.74)(2.10,0.89)};
  \foreach \px/\py in {0.30/0.50, 0.44/0.32, 0.58/0.42, 0.92/0.37,
                       1.08/0.19, 1.24/0.32, 1.54/0.61}
     {\fill[oiGreen!75!black] (\px,\py) circle (1.15pt);}
  \node[anchor=west, font=\bgfmtiny, text=oiGreen!45!black] at (1.14,0.10)
       {$e^{-E/kT}$};
\end{scope}

\node[tny, text width=3.46cm, anchor=north west] at (2.94,4.20)
  {$K$ perturbations of \emph{one} parent molecule\\
   (within-parent variance only)};

% ==================== BAND (c): local gradient =======================
\draw[bgfmpanel=oiVerm] (6.86,\By) rectangle (13.91,\BY);
\node[ttl, text=oiVerm!60!black] at (7.00,5.76) {(c) Local gradient};
\node[bdy, anchor=north west, text width=5.20cm] at (7.00,5.58)
     {closed-form score of the FM marginal\\[3pt]
      $s_\theta(x_t,t)=\dfrac{t\,v_\theta(x_t,t)-x_t}{(1-t)\,\sigma^{2}}$};

% mini inset: the two vectors that L_force aligns
\begin{scope}[shift={(12.30,5.30)}]
  \draw[oiGreyL, line width=0.4pt] (0,0) circle (0.44);
  \fill[oiGrey!85] (0,0) circle (1.7pt);
  \draw[-stealth', line width=0.9pt, draw=oiBlue] (0,0) -- (0.44,0.21);
  \draw[-stealth', line width=0.9pt, draw=oiVerm,
        dash pattern=on 2pt off 1.4pt] (0,0) -- (0.40,0.37);
  \draw[oiGrey, line width=0.35pt] (0.22,0.105) arc (25.6:42.4:0.245);
  \node[anchor=west, font=\bgfmtiny, text=oiBlue!75!black] at (0.46,0.14) {$s_\theta$};
  \node[anchor=west, font=\bgfmtiny, text=oiVerm!70!black] at (0.42,0.44) {$F/kT$};
\end{scope}

\node[lossbadge=oiVerm, text width=6.62cm, anchor=north west] at (7.00,4.60)
  {$\mathcal{L}_{\mathrm{force}}=\big\|\,s_\theta(x_{t^\star}\!,t^\star)
     -F(x_1)/kT\,\big\|_{2}^{2}$\\[1pt]
   {\bgfmtiny at the Boltzmann optimum
    $\nabla\!\log p=-\nabla E/kT=F/kT$;
    \textbf{measured harmful, \Cref{sec:negative}; $\lambda_1\!=\!0$ in every reported model}}};
\node[tny, text width=6.62cm, anchor=north west] at (7.00,3.80)
  {Read out at $t^{\star}\in\{0.85,0.92,0.97\}$ and averaged: the score has a
   $(1-t)^{-1}$ pole at $t\!=\!1$, so the per-atom norm is capped.};

% ==================== BAND (d): the physics bus ======================
\draw[physarr] (3.30,2.96) -- (3.30,3.34);
\node[anchor=west, font=\bgfmtiny, text=oiOrange!45!black] at (3.38,3.16)
     {$E(x_1^{(k)})/kT$};
\draw[physarr] (10.20,2.96) -- (10.20,3.34);
\node[anchor=west, font=\bgfmtiny, text=oiOrange!45!black] at (10.28,3.16)
     {$F(x_1)/kT=-\nabla E/kT$};

\draw[rounded corners=2.2pt, draw=oiOrange!85!black, line width=0.75pt,
      fill=oiOrange!12] (0.04,2.10) rectangle (13.91,2.92);
\node[anchor=west, font=\scriptsize\bfseries, text=oiOrange!35!black]
     at (0.18,2.66)
     {(d) Universal neural potential $E(x)$ --- OMol25 / eSEN};
\node[anchor=west, font=\bgfmtiny, text=oiGreyD] at (0.18,2.31)
     {3.9M molecules, 205M atoms, 83 elements; DFT forces \emph{and}
      energies; \textbf{no bond labels} --- one oracle, both terms};

% ==================== BAND (e): total loss + legend ==================
\draw[rounded corners=2pt, draw=oiGrey, line width=0.55pt, fill=white]
      (2.30,1.06) rectangle (11.66,1.82);
\node[anchor=center, font=\small] at (6.98,1.44)
  {$\mathcal{L}_{\mathrm{BGFM}}=\mathcal{L}_{\mathrm{FM}}
    +\lambda_1\,\textcolor{oiVerm}{\mathcal{L}_{\mathrm{force}}}
    +\lambda_2\,\textcolor{oiGreen}{\mathcal{L}_{\mathrm{energy}}}
    \qquad
    \textcolor{oiGreen}{\mathcal{L}_{\mathrm{energy}}}\!=\!0
    \iff p_\theta(\cdot\mid c)=\pi(\cdot\mid c)$};

\begin{scope}[shift={(0.30,0.28)}]
  \def\lw{0.26}
  \foreach \i/\c/\lbl in {%
      0/oiBlue/{model, flow matching},%
      1/oiGreen/{$\mathcal{L}_{\mathrm{energy}}$: global shape},%
      2/oiVerm/{$\mathcal{L}_{\mathrm{force}}$: local gradient},%
      3/oiOrange/{physics oracle (DFT labels)}}%
  {\pgfmathsetmacro{\xo}{\i*3.30}%
   \fill[\c!22] (\xo,0) rectangle (\xo+\lw,0.24);%
   \draw[\c, line width=0.6pt] (\xo,0) rectangle (\xo+\lw,0.24);%
   \node[anchor=west, font=\bgfmtiny] at (\xo+\lw+0.10,0.12) {\lbl};}
\end{scope}

\end{tikzpicture}

\bgfmFinish

%           \caption{...}\label{fig:method}
%         \end{figure}
%
%  SIZE.  Natural size 13.95cm x 8.70cm == 5.49in x 3.43in.
%  ICLR is a SINGLE 5.5in-wide text column (10pt Times, 11pt leading), so
%  this is a full-text-width figure at scale 1.0.  DO NOT rescale it: every
%  glyph inside is \footnotesize / \scriptsize / \tiny of the 10pt body
%  font, and scaling desynchronises in-figure text from the caption.
%  If a shorter figure is needed, delete band (e) (the total-loss box and
%  legend, y < 1.7) and state the loss in the caption instead -- that saves
%  1.5cm without touching any other coordinate.
%
%  LAYOUT (five bands, top to bottom).  Wire routing is deliberately
%  crossing-free:
%    (a) generative spine   prior noise -> flow matching -> x_1
%    (b) global-shape term  FFJORD log p  vs  E/kT     -> L_energy  [LEFT]
%    (c) local-gradient term closed-form score vs F/kT -> L_force   [RIGHT]
%    (d) physics bus        universal neural potential (OMol25 / eSEN)
%    (e) total loss + legend
%  (b) is on the left because the reverse-time (FFJORD) arrow terminates at
%  t=0 on the left; (c) is on the right because the score is read out at
%  t*=0.95 on the right; (d) is a full-width bus underneath so both terms
%  are fed by short vertical arrows and no wire ever crosses another.
% =====================================================================

% ---- standalone bootstrap (skipped when \input into a tikz-enabled doc) ----
\ifx\tikzpicture\undefined
  \documentclass[10pt]{article}
  \usepackage[paperwidth=14.30cm,paperheight=9.15cm,margin=0.15cm,noheadfoot]{geometry}
  \usepackage{amsmath,amssymb,bm}
  \usepackage{tikz}
  \setlength{\topskip}{0pt}
  \setlength{\parindent}{0pt}
  \pagestyle{empty}
  \begin{document}%
  \def\bgfmFinish{\end{document}}
\fi
\providecommand{\bgfmFinish}{\relax}
% \pending{...}: red placeholder marker, shared with the main text so that a
% single grep finds every not-yet-real number in the manuscript.
\providecommand{\pending}[1]{\textcolor{red}{#1}}
\providecommand{\bgfmFigScale}{1.0}

% 6pt/6.6pt: one step above \tiny (5pt at 10pt base), which is below the
% conventional 6pt floor for in-figure text.
\providecommand{\bgfmtiny}{\fontsize{6}{6.6}\selectfont}
\usetikzlibrary{positioning,calc,fit,backgrounds,decorations.pathmorphing,%
                arrows,shapes.geometric}

% ---------------------------------------------------------------- palette
% Okabe & Ito (2008): the de-facto standard CVD-safe qualitative set.
\definecolor{oiOrange}{HTML}{E69F00}
\definecolor{oiSky}   {HTML}{56B4E9}
\definecolor{oiGreen} {HTML}{009E73}
\definecolor{oiYellow}{HTML}{F0E442}
\definecolor{oiBlue}  {HTML}{0072B2}
\definecolor{oiVerm}  {HTML}{D55E00}
\definecolor{oiPurple}{HTML}{CC79A7}
\definecolor{oiGrey}  {HTML}{7F7F7F}
\definecolor{oiGreyL} {HTML}{BFBFBF}
\definecolor{oiGreyD} {HTML}{555555}
%
% SEMANTIC assignment -- identical in this figure and in the matplotlib
% figures (make_experiment_figures.py::PALETTE), so the paper reads as one
% visual system:
%   oiBlue    model / flow matching / learned quantities
%   oiGreen   L_energy  (global shape, FFJORD, Boltzmann)
%   oiVerm    L_force   (local gradient, score--force matching)
%   oiOrange  physics oracle (universal neural potential, DFT labels)
%   oiPurple  negative control (Y-scrambled)
%   oiGrey    no-physics baseline / neutral structure

% ---------------------------------------------------------------- styles
\tikzset{
  bgfmpanel/.style={rounded corners=2.2pt, draw=#1, line width=0.65pt,
                    fill=#1!5},
  ttl/.style={anchor=west, font=\footnotesize\bfseries, inner sep=0pt},
  bdy/.style={anchor=west, font=\scriptsize, inner sep=0pt, align=left},
  tny/.style={anchor=west, font=\bgfmtiny, inner sep=0pt, align=left,
              text=oiGreyD},
  flowarr/.style={-stealth', line width=0.9pt, draw=oiBlue},
  revarr/.style={-stealth', line width=0.75pt, draw=oiGreen,
                 dash pattern=on 2.2pt off 1.6pt},
  physarr/.style={-stealth', line width=0.8pt, draw=oiOrange!85!black},
  tapearr/.style={-stealth', line width=0.8pt, draw=oiBlue},
  lossbadge/.style={rounded corners=2pt, draw=#1, line width=0.6pt,
                    fill=#1!12, inner sep=2.6pt, align=left,
                    font=\scriptsize},
}

% ------------------------------------------------------- molecule glyphs
% Deterministic coordinates (no \pgfmathrandom) so the figure is
% byte-reproducible across compiles and across TeX distributions.
\newcommand{\bgfmAtom}[3]{\fill[#3] (#1,#2) circle (1.45pt);%
  \draw[oiGrey!60, line width=0.15pt] (#1,#2) circle (1.45pt);}
\newcommand{\bgfmAtomS}[3]{\fill[#3] (#1,#2) circle (0.95pt);%
  \draw[oiGrey!60, line width=0.15pt] (#1,#2) circle (0.95pt);}
\newcommand{\bgfmBond}[4]{\draw[oiGrey!85, line width=0.55pt] (#1,#2) -- (#3,#4);}

% t ~ 0 : unstructured point cloud, no bonds
\newcommand{\bgfmMolNoise}{%
  \bgfmAtomS{-0.34}{ 0.30}{oiGreyL}\bgfmAtomS{ 0.10}{ 0.36}{oiGreyL}%
  \bgfmAtomS{ 0.36}{ 0.06}{oiGreyL}\bgfmAtomS{-0.10}{-0.02}{oiGreyL}%
  \bgfmAtomS{-0.40}{-0.26}{oiGreyL}\bgfmAtomS{ 0.06}{-0.34}{oiGreyL}%
  \bgfmAtomS{ 0.40}{-0.32}{oiGreyL}\bgfmAtomS{ 0.24}{ 0.22}{oiGreyL}%
  \bgfmAtomS{-0.18}{ 0.12}{oiGreyL}}

% intermediate t : partially condensed, proto-structure
\newcommand{\bgfmMolMid}{%
  \bgfmBond{-0.20}{0.14}{0.04}{0.22}\bgfmBond{0.04}{0.22}{0.24}{0.02}%
  \bgfmBond{-0.20}{0.14}{-0.14}{-0.16}%
  \bgfmAtomS{-0.20}{ 0.14}{oiGrey!70}\bgfmAtomS{ 0.04}{ 0.22}{oiGrey!70}%
  \bgfmAtomS{ 0.24}{ 0.02}{oiGrey!70}\bgfmAtomS{-0.14}{-0.16}{oiGrey!70}%
  \bgfmAtomS{ 0.30}{ 0.30}{oiGreyL}\bgfmAtomS{-0.36}{-0.24}{oiGreyL}%
  \bgfmAtomS{ 0.20}{-0.30}{oiGreyL}}

% t = 1 : a clean structure -- 5-ring + N + O + a transition metal (green),
% so the "83 elements / transition metals" claim is visible in the glyph.
\newcommand{\bgfmMolFinal}{%
  \bgfmBond{0.000}{0.280}{-0.266}{0.087}\bgfmBond{-0.266}{0.087}{-0.165}{-0.227}%
  \bgfmBond{-0.165}{-0.227}{0.165}{-0.227}\bgfmBond{0.165}{-0.227}{0.266}{0.087}%
  \bgfmBond{0.266}{0.087}{0.000}{0.280}%
  \bgfmBond{0.266}{0.087}{0.560}{0.250}%
  \bgfmBond{-0.266}{0.087}{-0.520}{0.300}%
  \bgfmBond{-0.165}{-0.227}{-0.360}{-0.440}%
  \bgfmAtom{ 0.000}{ 0.280}{oiGrey!85}\bgfmAtom{-0.266}{ 0.087}{oiGrey!85}%
  \bgfmAtom{-0.165}{-0.227}{oiGrey!85}\bgfmAtom{ 0.165}{-0.227}{oiGrey!85}%
  \bgfmAtom{ 0.266}{ 0.087}{oiBlue}%
  \fill[oiGreen] (0.560,0.250) circle (2.4pt);%
  \draw[oiGrey!70, line width=0.2pt] (0.560,0.250) circle (2.4pt);%
  \bgfmAtomS{-0.520}{ 0.300}{oiVerm}%
  \bgfmAtomS{-0.360}{-0.440}{white}}

% =====================================================================
\begin{tikzpicture}[x=1cm, y=1cm, scale=\bgfmFigScale, >=stealth']
% guarantee a stable bounding box
\path (0,0.06) rectangle (13.95,8.66);

% ---- shared geometry -------------------------------------------------
\def\Ay{6.42}\def\AY{8.52}          % band (a) vertical extent
\def\bandL{2.30}\def\bandR{10.34}   % flow band horizontal extent
\def\tz{2.55}\def\to{10.09}         % x of t=0 and t=1 on the time axis
\def\tstar{9.71}                    % x of t* = 0.95
\def\By{3.36}\def\BY{5.96}          % bands (b)/(c) vertical extent

% ==================== BAND (a): generative spine =====================
% -- prior --------------------------------------------------------------
\draw[bgfmpanel=oiBlue] (0.04,\Ay) rectangle (2.02,\AY);
\node[anchor=north west, font=\scriptsize\bfseries] at (0.12,8.48) {(a)};
\begin{scope}[shift={(1.06,7.88)}]\bgfmMolNoise\end{scope}
\node[anchor=north, font=\scriptsize] at (1.03,7.44)
     {$x_0\!\sim\!\mathcal{N}(0,\sigma^2 I)$};
\node[anchor=north, font=\bgfmtiny, align=center, text=oiGreyD]
      at (1.03,7.02) {CTMC prior for\\types, charges};

% -- flow band ----------------------------------------------------------
\shade[left color=oiBlue!7, right color=oiGreen!9, rounded corners=2.2pt]
      (\bandL,\Ay) rectangle (\bandR,\AY);
\draw[rounded corners=2.2pt, draw=oiBlue, line width=0.65pt]
      (\bandL,\Ay) rectangle (\bandR,\AY);

% forward (generative) ODE
\draw[flowarr] (\tz,8.02) -- (\to,8.02);
\node[anchor=south, font=\scriptsize, text=oiBlue!75!black] at (6.32,8.08)
     {$\dfrac{d x_t}{dt}=v_\theta(x_t,t)$\;\; SE(3)-equivariant flow matching};

% the three glyphs along the path
\begin{scope}[shift={(3.40,7.40)}]\bgfmMolNoise\end{scope}
\begin{scope}[shift={(6.32,7.40)}]\bgfmMolMid\end{scope}
\begin{scope}[shift={(9.10,7.40)}]\bgfmMolFinal\end{scope}

% reverse-time (FFJORD) ODE: right -> left, terminates at t = 0
\draw[revarr] (\to,6.98) -- (\tz,6.98);
\node[anchor=north, font=\bgfmtiny, text=oiGreen!50!black] at (6.32,6.94)
     {FFJORD reverse-time ODE $+$ Hutchinson divergence};

% time axis
\draw[oiGrey, line width=0.5pt] (\tz,6.66) -- (\to,6.66);
\foreach \xx in {\tz,\tstar,\to}%
  {\draw[oiGrey, line width=0.5pt] (\xx,6.66) -- (\xx,6.58);}
\node[anchor=north, font=\bgfmtiny] at (\tz,6.57) {$t\!=\!0$};
\node[anchor=north east, font=\bgfmtiny] at (9.60,6.57) {$t^{\star}\!<\!1$};
\node[anchor=north west, font=\bgfmtiny] at (10.13,6.57) {$1$};

% -- x_1 ----------------------------------------------------------------
\draw[bgfmpanel=oiBlue] (10.62,\Ay) rectangle (13.91,\AY);
\node[anchor=north, font=\scriptsize] at (12.26,8.46) {$x_1$};
\node[anchor=north, font=\bgfmtiny, align=center] at (12.26,8.04)
     {positions $+$ atomic numbers\\\textbf{83 elements, bond-free}};
\node[anchor=north, font=\bgfmtiny, align=center, text=oiGreyD]
     at (12.26,7.10) {bonds recovered \emph{post hoc}\\(xyz2mol), never supervised};

% ==================== taps into the two branches =====================
\draw[revarr]  (\tz,\Ay)    -- (\tz,\BY-0.02);
\node[anchor=west, font=\bgfmtiny, text=oiGreen!50!black] at (\tz+0.09,6.20)
     {$\log p_\theta(x_1)$};
\draw[tapearr] (\tstar,\Ay) -- (\tstar,\BY-0.02);
\node[anchor=west, font=\bgfmtiny, text=oiBlue!75!black] at (\tstar+0.09,6.20)
     {$v_\theta(x_{t^\star},t^\star)$};

% ==================== BAND (b): global shape =========================
\draw[bgfmpanel=oiGreen] (0.04,\By) rectangle (6.58,\BY);
\node[ttl, text=oiGreen!45!black] at (0.18,5.76)
     {(b) Global shape: energy consistency};
\node[bdy, anchor=north west] at (0.18,5.58)
     {$\log p_\theta(x_1)=\log p_0(x_0)-\displaystyle\int_{0}^{1}\!\!
       \nabla\!\cdot\!v_\theta\,dt$};

% full-panel-width loss badge: the formula stays on one line
\node[lossbadge=oiGreen, text width=6.06cm, anchor=north west] at (0.18,5.02)
  {$\mathcal{L}_{\mathrm{energy}}=
    \operatorname{Var}_{k}\!\big[\log p_\theta(x_1^{(k)})+E(x_1^{(k)})/kT\big]$\\[1pt]
   {\bgfmtiny the partition function $Z$ cancels exactly}};

% mini potential-energy surface + the K within-parent perturbations
\begin{scope}[shift={(0.30,3.46)}]
  \draw[oiGrey, line width=0.45pt] (0,0) -- (2.26,0);
  \draw[oiGrey, line width=0.45pt] (0,0) -- (0,0.90);
  \node[font=\bgfmtiny, rotate=90, inner sep=1pt] at (-0.14,0.45) {$E$};
  % Boltzmann weight e^{-E/kT} shaded under the curve
  \fill[oiGreen!16]
    plot[smooth] coordinates {(0.12,0.09)(0.44,0.53)(0.76,0.17)(1.08,0.64)
                              (1.42,0.26)(1.76,0.12)(2.10,0.06)}
    -- (2.10,0) -- (0.12,0) -- cycle;
  % E(x): a double well
  \draw[oiBlue, line width=0.8pt]
    plot[smooth] coordinates {(0.12,0.84)(0.44,0.32)(0.76,0.66)(1.08,0.19)
                              (1.42,0.54)(1.76,0.74)(2.10,0.89)};
  \foreach \px/\py in {0.30/0.50, 0.44/0.32, 0.58/0.42, 0.92/0.37,
                       1.08/0.19, 1.24/0.32, 1.54/0.61}
     {\fill[oiGreen!75!black] (\px,\py) circle (1.15pt);}
  \node[anchor=west, font=\bgfmtiny, text=oiGreen!45!black] at (1.14,0.10)
       {$e^{-E/kT}$};
\end{scope}

\node[tny, text width=3.46cm, anchor=north west] at (2.94,4.20)
  {$K$ perturbations of \emph{one} parent molecule\\
   (within-parent variance only)};

% ==================== BAND (c): local gradient =======================
\draw[bgfmpanel=oiVerm] (6.86,\By) rectangle (13.91,\BY);
\node[ttl, text=oiVerm!60!black] at (7.00,5.76) {(c) Local gradient};
\node[bdy, anchor=north west, text width=5.20cm] at (7.00,5.58)
     {closed-form score of the FM marginal\\[3pt]
      $s_\theta(x_t,t)=\dfrac{t\,v_\theta(x_t,t)-x_t}{(1-t)\,\sigma^{2}}$};

% mini inset: the two vectors that L_force aligns
\begin{scope}[shift={(12.30,5.30)}]
  \draw[oiGreyL, line width=0.4pt] (0,0) circle (0.44);
  \fill[oiGrey!85] (0,0) circle (1.7pt);
  \draw[-stealth', line width=0.9pt, draw=oiBlue] (0,0) -- (0.44,0.21);
  \draw[-stealth', line width=0.9pt, draw=oiVerm,
        dash pattern=on 2pt off 1.4pt] (0,0) -- (0.40,0.37);
  \draw[oiGrey, line width=0.35pt] (0.22,0.105) arc (25.6:42.4:0.245);
  \node[anchor=west, font=\bgfmtiny, text=oiBlue!75!black] at (0.46,0.14) {$s_\theta$};
  \node[anchor=west, font=\bgfmtiny, text=oiVerm!70!black] at (0.42,0.44) {$F/kT$};
\end{scope}

\node[lossbadge=oiVerm, text width=6.62cm, anchor=north west] at (7.00,4.60)
  {$\mathcal{L}_{\mathrm{force}}=\big\|\,s_\theta(x_{t^\star}\!,t^\star)
     -F(x_1)/kT\,\big\|_{2}^{2}$\\[1pt]
   {\bgfmtiny at the Boltzmann optimum
    $\nabla\!\log p=-\nabla E/kT=F/kT$;
    \textbf{measured harmful, \Cref{sec:negative}; $\lambda_1\!=\!0$ in every reported model}}};
\node[tny, text width=6.62cm, anchor=north west] at (7.00,3.80)
  {Read out at $t^{\star}\in\{0.85,0.92,0.97\}$ and averaged: the score has a
   $(1-t)^{-1}$ pole at $t\!=\!1$, so the per-atom norm is capped.};

% ==================== BAND (d): the physics bus ======================
\draw[physarr] (3.30,2.96) -- (3.30,3.34);
\node[anchor=west, font=\bgfmtiny, text=oiOrange!45!black] at (3.38,3.16)
     {$E(x_1^{(k)})/kT$};
\draw[physarr] (10.20,2.96) -- (10.20,3.34);
\node[anchor=west, font=\bgfmtiny, text=oiOrange!45!black] at (10.28,3.16)
     {$F(x_1)/kT=-\nabla E/kT$};

\draw[rounded corners=2.2pt, draw=oiOrange!85!black, line width=0.75pt,
      fill=oiOrange!12] (0.04,2.10) rectangle (13.91,2.92);
\node[anchor=west, font=\scriptsize\bfseries, text=oiOrange!35!black]
     at (0.18,2.66)
     {(d) Universal neural potential $E(x)$ --- OMol25 / eSEN};
\node[anchor=west, font=\bgfmtiny, text=oiGreyD] at (0.18,2.31)
     {3.9M molecules, 205M atoms, 83 elements; DFT forces \emph{and}
      energies; \textbf{no bond labels} --- one oracle, both terms};

% ==================== BAND (e): total loss + legend ==================
\draw[rounded corners=2pt, draw=oiGrey, line width=0.55pt, fill=white]
      (2.30,1.06) rectangle (11.66,1.82);
\node[anchor=center, font=\small] at (6.98,1.44)
  {$\mathcal{L}_{\mathrm{BGFM}}=\mathcal{L}_{\mathrm{FM}}
    +\lambda_1\,\textcolor{oiVerm}{\mathcal{L}_{\mathrm{force}}}
    +\lambda_2\,\textcolor{oiGreen}{\mathcal{L}_{\mathrm{energy}}}
    \qquad
    \textcolor{oiGreen}{\mathcal{L}_{\mathrm{energy}}}\!=\!0
    \iff p_\theta(\cdot\mid c)=\pi(\cdot\mid c)$};

\begin{scope}[shift={(0.30,0.28)}]
  \def\lw{0.26}
  \foreach \i/\c/\lbl in {%
      0/oiBlue/{model, flow matching},%
      1/oiGreen/{$\mathcal{L}_{\mathrm{energy}}$: global shape},%
      2/oiVerm/{$\mathcal{L}_{\mathrm{force}}$: local gradient},%
      3/oiOrange/{physics oracle (DFT labels)}}%
  {\pgfmathsetmacro{\xo}{\i*3.30}%
   \fill[\c!22] (\xo,0) rectangle (\xo+\lw,0.24);%
   \draw[\c, line width=0.6pt] (\xo,0) rectangle (\xo+\lw,0.24);%
   \node[anchor=west, font=\bgfmtiny] at (\xo+\lw+0.10,0.12) {\lbl};}
\end{scope}

\end{tikzpicture}

\bgfmFinish
